
\section{Introduction}
\label{sec:introduction}

\subsection{The Simulation Gap in Silicon Spin Qubit Control}
\label{subsec:simulation_gap}

Silicon spin qubits—particularly SiMOS devices—are poised to scale quantum computing. With coherence times approaching $T_1 \sim 1$\,ms and $T_2 \sim 0.5$\,ms~\cite{Xue2022,Noiri2022,Philips2022,Steinacker2023}, CMOS fabrication compatibility, and natural nearest-neighbour connectivity, they offer a manufacturable path toward fault-tolerant architectures. Yet their noise sensitivity remains a critical roadblock: charge fluctuations, hyperfine coupling variations, and residual spin-orbit interactions degrade gate fidelity in ways that first-principles prediction cannot reliably capture~\cite{Veldhorst2015,Fogarty2018}.

Closing this gap demands simulation environments that are simultaneously physically rigorous and machine-learning accessible. Current tools fall short on at least one dimension. Lindblad solvers such as QuTiP~\cite{Lambert2024} deliver precise open-system dynamics but provide no reinforcement learning interface. Quantum RL environments~\cite{VanDerLinde2023,QiskitGym2024} offer Gymnasium compatibility but rely on abstract depolarising noise rather than device-specific physics. SiliQun resolves this divide as the first platform to merge device-accurate Lindblad simulation with a standardised RL interface for silicon spin qubits.

\subsection{The Physics-Grounded RL Simulation Methodology}
\label{subsec:intro_pgirs}

We introduce the \textbf{Physics-Grounded RL Simulation} (PGIRS) methodology for building platforms that are both physically rigorous and algorithmically accessible. PGIRS rests on five design principles:

\begin{itemize}
  \item \textbf{Noise-First Design:} Noise models drive all simulator components, validated against physical benchmarks before RL integration.
  \item \textbf{Gymnasium Compliance:} A standard Gymnasium \texttt{Env} interface ensures plug-and-play compatibility with RL libraries.
  \item \textbf{Reward Engineering:} Bayesian optimisation automatically discovers reward weights, eliminating manual tuning.
  \item \textbf{Hierarchical Decomposition:} Multi-qubit objectives decompose into sub-goals with dense intermediate rewards.
  \item \textbf{Empirical Validation:} Each component undergoes analytical benchmarking, cross-algorithm testing, and hardware transfer checks.
\end{itemize}

SiliQun implements PGIRS across the silicon spin qubit family. Tool limitations in the field, however, have so far prevented adoption of this methodology at scale.

\subsection{Limitations of Existing Frameworks}
\label{subsec:intro_limitations}

No open-source framework currently satisfies both requirements—physical accuracy and standard RL support—for silicon spin qubit simulation. QuTiP~\cite{Lambert2024} offers high-fidelity Lindblad solutions but lacks Gymnasium integration. Qiskit Pulse~\cite{Alexander2020} enables pulse-level control for superconducting qubits but omits SiMOS models entirely. C3~\cite{Wittler2021} supports closed-loop optimisation via GRAPE and CMA-ES yet excludes standard RL algorithms. TensorFlow Quantum~\cite{Broughton2020} targets gate-model circuits via Cirq, not spin qubit physics. Quantum RL environments (qgym~\cite{VanDerLinde2023}, Qiskit Gym~\cite{QiskitGym2024}) use simplified depolarising noise, making them unsuitable for SiMOS-specific dynamics. Orbell's Oxford work~\cite{Orbell2024} advances Bayesian device tuning for SiMOS but operates at device characterisation rather than gate synthesis, and provides no Gymnasium interface. Unlike these tools, SiliQun's SDQF architecture encodes hardware-specific noise profiles—SiMOS charge noise, donor hyperfine coupling, GAA spin-orbit interaction—into a reusable RL environment, enabling direct policy training on device-realistic dynamics.

\subsection{Contributions}
\label{subsec:contributions}

This work advances three interconnected areas, ordered by centrality to SiliQun:

\textbf{Primary contributions --- the SiliQun simulation platform:}
\begin{enumerate}
  \item \textbf{SiliQun physics engine:} A unified Lindblad solver supporting SiMOS, donor-electron, and GAA qubits via configurable hardware profiles, validated at numerical precision ($\max\,\delta \leq 2.78 \times 10^{-16}$) and GPU-accelerated for $93\times$ faster training cycles.
  \item \textbf{Multi-technology hardware profiles:} A \texttt{TechnologyProfile} abstraction encoding device physics for SiMOS, donor-electron, GAA, GaAs singlet-triplet, and SLEDGE spin qubits, grounded in published experimental data and selectable via a single constructor argument.
  \item \textbf{Gymnasium-compatible interface:} A discrete 11-token gate action space, 16-element density-matrix state representation, and auto-weighted four-component reward function in a fully compliant \texttt{SiliQunEnv}.
  \item \textbf{Qiskit BackendV2 provider:} The pip-installable \texttt{siliqun-provider} exposes SiliQun as a Qiskit \texttt{BackendV2}-compatible backend, integrating directly with the Qiskit ecosystem.
\end{enumerate}
